\documentclass[11pt]{article}
\usepackage[T1]{fontenc}
\usepackage[utf8]{inputenc}
\usepackage[margin=1in]{geometry}
\usepackage{amsmath,amssymb,amsthm,microtype}
\usepackage[colorlinks=true,linkcolor=blue,citecolor=blue,urlcolor=blue]{hyperref}
\newtheorem{theorem}{Theorem}
\newtheorem{proposition}[theorem]{Proposition}
\newcommand{\Z}{\mathbb Z}
\newcommand{\Pp}{\mathbb P}
\newcommand{\E}{\mathbb E}
\newcommand{\pc}{p_{\mathrm c}}
\newcommand{\bd}{\partial}
\title{Finite-energy lattice percolation with\break internal critical probability one}
\author{Alec Kriebel\\\small Independent researcher\\
\small\href{https://orcid.org/0009-0001-9320-500X}{ORCID: 0009-0001-9320-500X}}
\date{October 15, 2026\quad---\quad Version 1.0.0}
\begin{document}
\maketitle
\begin{abstract}
For every integer $d\ge2$ we construct a bond percolation on $\Z^d$ that is invariant under all lattice automorphisms, mixing under translations, and has ordinary finite energy. Almost surely it has a unique infinite component, yet every independent Bernoulli thinning with retention $p<1$ destroys all infinite components. The construction perturbs a one-ended spanning tree, using the maximum finite-past size along each non-tree edge's tree path to suppress every bypass of large tree cuts. This answers the historical ordinary-energy question of Benjamini, H\"aggstr\"om and Schramm. We also make explicit a planar threshold-one consequence of H\"aggstr\"om--Mester's earlier construction. No first-priority claim is made.
\end{abstract}

\section{Statement and context}
Write $E_d$ for the nearest-neighbor edges of $\Z^d$. A bond configuration $X\subseteq E_d$ has \emph{ordinary finite energy} if, for every edge $e$,
\begin{equation}\label{energy}
 0<\Pp(e\in X\mid\sigma(\mathbf1_{\{f\in X\}}:f\ne e))<1
 \quad\text{almost surely}.
\end{equation}
No common positive lower bound is required. For a fixed countable graph $H$, define $\pc(H)$ as the infimum of retention parameters at which independent bond thinning of $H$ has an infinite component with positive probability. The internal site threshold is defined analogously. These definitions apply to disconnected graphs as well.

Benjamini, H\"aggstr\"om and Schramm asked whether an automorphism-invariant, finite-energy bond percolation can percolate almost surely while satisfying $\pc(X)=1$ \cite[Question 1.1]{BHS}. Their finite-energy convention is \eqref{energy}. The question reappears in Benjamini's lecture notes \cite[Open problem 8.15, p.~61 of the notes version]{S} and Benjamini--Tassion \cite[Question 1]{BT}. Finite components are allowed: almost-sure percolation means existence of an infinite component.

\begin{theorem}\label{main}
For every integer $d\ge2$, there is a bond law on $(\Z^d,E_d)$ invariant under every graph automorphism and mixing under translations, satisfying \eqref{energy}, such that almost surely $X$ has a unique infinite component $C$ and
\[
 \pc(X)=\pc(C)=1.
\]
The internal site thresholds of $X$ and $C$ also equal one. All threshold assertions are quenched: the realized graph is fixed before introducing independent thinning.
\end{theorem}

The construction develops the summable finite-past perturbations of H\"aggstr\"om and Mester \cite[Section 2]{HM}, using established spanning-tree inputs of Pemantle and Tim\'ar \cite{P,T}. Its specific rule controls all non-tree bypasses in any bounded-degree graph with a one-ended spanning tree. Section~\ref{planar} shows that the planar existence conclusion also follows from an exact midpoint projection of the older HM model, with an additional bounded-cutset argument. That deduction is not stated in the inspected HM paper. The distinction precludes attributing a wholly new planar mechanism to this note, or claiming a verified first resolution of the historical question. The first occurrence of the present rule and all-dimensional conclusion is not certified by the accompanying bounded literature audit.

\section{A perturbation of a one-ended spanning tree}
Let $G=(V,E)$ be a countably infinite connected simple graph with degree at most $\Delta<\infty$, and let $T$ be a one-ended spanning tree. For $t\in T$, deleting $t$ leaves exactly one finite component $S_T(t)$; set $s_T(t)=|S_T(t)|\ge1$. For $e=\{x,y\}\notin T$, let $P_T(e)$ be the unique finite tree path between its endpoints and put
\[
 M_T(e)=\max_{t\in P_T(e)}s_T(t).
\]
Conditionally on $T$, choose the output edges independently with probabilities
\begin{equation}\label{kernel}
 r_T(e)=
 \begin{cases}
 1-2^{-s_T(e)},&e\in T,\\
 2^{-M_T(e)},&e\notin T.
 \end{cases}
\end{equation}
All these parameters lie strictly between zero and one.

\begin{proposition}\label{kernelprop}
For any possibly random one-ended spanning-tree input $T$, the output of \eqref{kernel} has ordinary finite energy. Almost surely it has a unique infinite component, and its internal bond and site thresholds, and those of that component, equal one. The kernel is measurable and commutes with every automorphism of $G$.
\end{proposition}

\begin{proof}
Fix a good deterministic input tree. From every $v\in V$ there is a unique infinite tree ray, with successive edges $t_0,t_1,\ldots$. Write $S_n=S_T(t_n)$ and $s_n=|S_n|$. The $S_n$ contain $v$ and increase strictly: the next ray vertex belongs to $S_{n+1}\setminus S_n$. Thus the $s_n$ are distinct positive integers. Rays in a one-ended tree coalesce, and consequently the $S_n$ exhaust $V$.

The only tree edge crossing $S_n$ is $t_n$. If a non-tree edge $e$ crosses $S_n$, its tree path uses $t_n$, so $M_T(e)\ge s_n$. There are at most $\Delta s_n$ ambient crossing edges. Let $B_n$ be the event that $t_n$ is deleted or some non-tree crossing edge is inserted. The union bound gives
\begin{equation}\label{sum}
 \sum_n\Pp(B_n\mid T)
 \le \sum_{m\ge1}(1+\Delta m)2^{-m}
 =1+2\Delta<\infty.
\end{equation}
The first Borel--Cantelli lemma requires no independence between the $B_n$. It yields, eventually along the ray,
\begin{equation}\label{singleton}
 t_n\in X,\qquad \bd_X S_n=\{t_n\}.
\end{equation}
Intersect over countably many starting vertices and then integrate over $T$. We obtain one joint probability-one event on which \eqref{singleton} holds eventually at every vertex.

On that event the retained tree $R=X\cap T$ contains infinite ray tails. Their coalescence shows that $R$ has a unique infinite component $C_R$: any other infinite component would contain a tree ray and hence meet a common retained tail. Every infinite $X$-component starting at $v$ must leave each finite $S_n$. At a sufficiently late cut it must cross $t_n$, both of whose endpoints belong to $C_R$. Thus every infinite $X$-component meets $C_R$, proving uniqueness of the infinite component $C$ of $X$.

Now fix a good pair $(T,X)$ and introduce independent bond thinning with retention $p<1$. An infinite path from $v$ must retain every sufficiently late $t_n$. These are distinct fixed thinning coordinates, so the probability of such a path is at most $p^k$ for every $k$, and is zero. A countable union over vertices excludes all infinite thinned components. The same deterministic witness works for every real $p<1$, without intersecting an uncountable collection of probability-one events. At $p=1$ the graph percolates. For $C$, the late cut identity is
\[
 \bd_C(S_n\cap C)=\{t_n\},
\]
since both endpoints lie in $C_R\subseteq C$. This proves both bond thresholds. An infinite site-thinned path must retain infinitely many distinct vertices of the same ray, giving the identical $p^k$ obstruction and both site thresholds.

For completeness, the witness-to-output step is measurable. At fixed rational $p$ and vertex $v$, the thinning probability of an infinite connection is a decreasing limit of measurable finite-ball connection probabilities. Requiring these limits to vanish for countably many $p,v$ is Borel; monotonicity supplies all real $p<1$. Alternatively, the bond obstruction has the graph-only Borel certificate that, for every $v$ and $k$, there are finite sets containing $v$ with singleton boundaries at $k$ pairwise distinct edges. No measurable selection of a hidden tree from $X$ is needed.

To prove energy after forgetting $T$, realize \eqref{kernel} by independent uniforms $U_e$, independent of $T$. Put $\mathcal F_e=\sigma(X_f:f\ne e)$ and $\mathcal H_e=\sigma(T,U_f:f\ne e)$. Then $\mathcal F_e\subseteq\mathcal H_e$, $U_e$ is independent of $\mathcal H_e$, and
\begin{equation}\label{tower}
 \Pp(X_e=1\mid\mathcal F_e)=\E[r_T(e)\mid\mathcal F_e],\qquad
 \Pp(X_e=0\mid\mathcal F_e)=\E[1-r_T(e)\mid\mathcal F_e].
\end{equation}
Both conditional expectations are strictly positive almost surely. Indeed a conditional expectation of a strictly positive integrable random variable cannot vanish on a set of positive probability, by integration over that set. Countability of edges gives a common null set. Conditional versions may be assigned value $1/2$ on excluded conditioning configurations.

Finally, the finite-side and path functions are Borel: enumerate finite connected vertex sets with their finite incident boundaries and finite paths. Connectedness, acyclicity and one-endedness of a spanning tree also have countable Borel descriptions. Assigning rates $1/2$ outside the probability-one good input domain gives a globally measurable kernel. Finite-side cardinalities and tree paths commute with automorphisms, so the kernel does too.
\end{proof}

\section{Lattice input and symmetry}
Tim\'ar's one-ended spanning-tree theorem \cite[Theorem 1 and Corollary 2]{T} supplies a one-ended spanning tree as a measurable isomorphism-equivariant factor of iid vertex labels on an ergodic amenable one-ended unimodular random graph. We use the precise inspected version arXiv:1805.10690v1. The deterministic rooted lattice $\Z^d$, $d\ge2$, meets these hypotheses. The mass-transport identity follows by summing over displacement vectors; invariant uniformly shifted finite box partitions verify amenability; and the complement of a sufficiently large box is connected, verifying one-endedness. Its rooted isomorphism class is deterministic and hence ergodic.

This factor input and independent iid edge marks have a joint law invariant under every lattice automorphism. Proposition~\ref{kernelprop} therefore gives the required full symmetry. The joint iid field is mixing under translations: finite-coordinate events are independent after sufficiently large translations, and approximation in probability by such events proves the assertion for arbitrary measurable events. An equivariant measurable factor inherits mixing. Proposition~\ref{kernelprop} now proves Theorem~\ref{main}.

For $2\le d\le4$, Pemantle's one-ended UST gives an older alternative input \cite[Theorems 2.3 and 4.3]{P}. In higher dimensions that UST limit is disconnected, so its substitution into \eqref{kernel} would leave some connecting tree paths undefined. Tim\'ar's connected input avoids this issue. No uniform energy bound, finite coding radius or quantitative mixing rate is asserted.

\section{The planar antecedent}\label{planar}
We record explicitly why the older HM construction deserves stronger credit than similarity of its flip rule. In the square lattice let $T$ be the primal UST and $T^*$ its complementary one-ended dual tree, as in \cite[Section 2]{HM}. The exact fixed-phase midpoint coordinates of HM's scale-two site model, interpreted as coarse bonds with all coarse vertices present, have conditional probabilities
\begin{equation}\label{hmkernel}
 q_T(e)=
 \begin{cases}
 1-4^{-s_T(e)},&e\in T,\\
 4^{-s_{T^*}(e^*)},&e\notin T.
 \end{cases}
\end{equation}
Indeed HM flips a site $x$ with probability $2^{-b(x)-1}$, where $b(x)$ counts the finite vertices separated from monochromatic infinity by removing $x$. A midpoint with a finite coarse side of $s$ vertices has $b(x)=s+(s-1)=2s-1$. Taking these coordinates preserves their conditional independence. This is a coordinate marginal of the fixed-phase construction, not the induced graph on its finally open site set. The random parity phase of the full site model is absent from the coarse bond law.

Here is the additional threshold argument. Along a finite simple path in a one-ended tree, every finite-side edge size occurs at most twice: the path is two increasing ancestor chains meeting at their common ancestor. For the ray cuts $S_n$ of Section 2, the duals of $\bd_G S_n$ form exactly the fundamental cycle obtained by adding $t_n^*$ to $T^*$. To see this directly, a finite primal cut is an even dual edge set; all its edges but $t_n^*$ lie in the dual tree, so its only possible cycle is the simple path in $T^*$ between the endpoints of $t_n^*$ together with that edge. No additional even component can lie in a tree.

Thus the non-tree edges crossing $S_n$ correspond to one finite dual-tree path. If $Y$ denotes the output of \eqref{hmkernel},
\[
 \E[|\bd_Y S_n|\mid T]
 \le 1+2\sum_{m\ge1}4^{-m}=\frac53.
\]
Conditional Fatou implies $\liminf_n|\bd_Y S_n|<\infty$ almost surely. Since $S_n$ exhaust the lattice, their boundaries escape every finite edge set. For each good output, choose infinitely many pairwise edge-disjoint boundaries of size at most some finite $K$. If a boundary is empty there is no infinite escape from its interior. Otherwise a Bernoulli-$p$ thinning deletes each entire boundary with probability at least $(1-p)^K$. Independence on the disjoint boundaries bounds the probability of crossing $k$ of them by $[1-(1-p)^K]^k\to0$. Every fixed vertex lies in all sufficiently late $S_n$, so this quenched argument excludes every infinite component for every $p<1$.

Conversely, summable deletions along primal rays leave infinite retained tails. The coarse law is fully automorphism invariant, and the conditional-expectation argument \eqref{tower} proves its ordinary finite energy. Consequently this exact HM midpoint projection already gives the historical bond existence conclusion in dimension two, once the bounded-cutset deduction is supplied. The actual post-flip HM site model also has internal site threshold one: its surrounding fine-lattice circuits have one initially open midpoint and a dual-tree path whose expected number of flipped-open sites is at most $2\sum_{m\ge0}2^{-m-1}=2$. Fatou and disjoint surrounding circuits give the same thinning obstruction. The internal bond threshold follows from \cite[Lemma 3.1]{HM}; that lemma by itself does not convert an external site law into a bond law with finite energy.

The inspected HM paper states finite-energy coexistence and discusses threshold one for its \emph{pre-flip} model; it does not state the post-flip deduction above. The inspected 2015 version of \cite{BT} still poses the exact bond question. We neither erase this chronology nor infer a previously printed solution from our deduction. The present rule \eqref{kernel} uses no planar dual tree and gives eventual singleton cuts, a stronger structural conclusion than the bounded-cutset subsequence used here. Its mathematical validity is independent of the unresolved historical-priority question.

\paragraph{Verification and disclosure.}
The accompanying archive contains editable source, the frozen original proof and reviews, independent analytical audits of the theorem and planar antecedent, a bounded priority audit with source/version locators, and reproducible byte-integrity checks. These checks do not certify the infinite mathematical proof or priority. Downloaded publications are not redistributed. OpenAI tools were used extensively in solving, proof development, drafting, and adversarial verification. This is an unrefereed preprint; no completed external human peer review or formal proof-assistant verification is claimed.

\begin{thebibliography}{9}
\bibitem{BHS} I.~Benjamini, O.~H\"aggstr\"om and O.~Schramm,
On the effect of adding $\varepsilon$-Bernoulli percolation to everywhere percolating subgraphs of $\Z^d$,
\emph{J. Math. Phys.} \textbf{41} (2000), 1294--1297.
\href{https://doi.org/10.1063/1.533187}{doi:10.1063/1.533187}; inspected arXiv:math/9906002v1 (1999).
\bibitem{S} I.~Benjamini, \emph{Coarse Geometry and Randomness}, Lecture Notes in Mathematics 2100, Springer, 2013.
The problem locator here refers to the author's lecture-notes version, Open problem 8.15, p.~61, in the \href{https://arquivo.pt/noFrame/replay/20201231041548id_/http://www.wisdom.weizmann.ac.il/~itai/stflouraug24.pdf}{archived PDF}.
\bibitem{BT} I.~Benjamini and V.~Tassion, Homogenization via sprinkling,
\emph{Ann. Inst. Henri Poincar\'e Probab. Stat.} \textbf{53} (2017), 997--1005.
\href{https://doi.org/10.1214/16-AIHP746}{doi:10.1214/16-AIHP746}; inspected arXiv:1505.06069v1 (2015).
\bibitem{HM} O.~H\"aggstr\"om and P.~Mester, Some two-dimensional finite energy percolation processes,
\emph{Electron. Commun. Probab.} \textbf{14} (2009), 42--54.
\href{https://doi.org/10.1214/ECP.v14-1446}{doi:10.1214/ECP.v14-1446}; inspected published PDF.
\bibitem{P} R.~Pemantle, Choosing a spanning tree for the integer lattice uniformly,
\emph{Ann. Probab.} \textbf{19} (1991), 1559--1574.
Inspected \href{https://arxiv.org/abs/math/0404043v1}{arXiv:math/0404043v1} (deposited 2004).
\bibitem{T} \'A.~Tim\'ar, One-ended spanning trees in amenable unimodular graphs,
\emph{Electron. Commun. Probab.} \textbf{24} (2019), paper 72, 12~pp.
\href{https://doi.org/10.1214/19-ECP274}{doi:10.1214/19-ECP274}; theorem used in the inspected 10-page \href{https://arxiv.org/abs/1805.10690v1}{arXiv:1805.10690v1} (2018).
\end{thebibliography}
\end{document}
